\documentclass[11pt]{article}

\usepackage[margin=1.1in]{geometry}
\usepackage{amsmath,amssymb,amsthm}
\usepackage{booktabs}
\usepackage{microtype}
\usepackage{xcolor}
\definecolor{linkblue}{RGB}{25,60,120}
\usepackage[colorlinks=true,linkcolor=linkblue,citecolor=linkblue,urlcolor=linkblue]{hyperref}

\newtheorem{definition}{Definition}
\newtheorem{remark}{Remark}

\newcommand{\F}{\mathbb{F}}
\newcommand{\Ftwo}{\F_2}
\newcommand{\Fsixteen}{\F_{16}}
\newcommand{\Ftag}{\F_{2^{128}}}
\newcommand{\shake}{\mathsf{SHAKE256}}
\newcommand{\ctx}{\mathsf{ctx}}
\newcommand{\Spre}{\mathsf{SPre}}
\newcommand{\enc}{\mathsf{enc}_{64}}
\newcommand{\Cred}{\mathsf{C}}
\newcommand{\Null}{\mathsf{N}}
\newcommand{\Def}{\mathsf{D}}

\title{VOLE-ACT: Post-Quantum Anonymous Credit Tokens\\
from VOLE-in-the-Head and MAYO}
\author{Samuel Schlesinger}
\date{July 2026\\[0.5em]\normalsize Research prototype documentation --- not yet proven or independently audited}

\begin{document}
\maketitle

\begin{abstract}
Anonymous credit tokens let an issuer grant a client a hidden, spendable
balance: the client can later redeem credits unlinkably, each redemption
consumes a one-time nullifier and returns a fresh token on the remaining
balance. Classical constructions rest on discrete-logarithm-style algebra ---
blind or randomizable signatures, algebraic MACs, Sigma protocols --- none of
which survives a quantum adversary, and the post-quantum blind-signature
literature has proved fragile. We describe VOLE-ACT, a post-quantum
instantiation of the Katz--Schlesinger anonymous-credit design whose entire
symmetric layer is Keccak and whose only algebraic assumption is the
multivariate-quadratic (whipped MAYO) trapdoor. The central idea is to
\emph{avoid blindness entirely}: the issuer signs a revealed SHAKE256
commitment digest whose preimage stays hidden, and possession of the signature
is later proven in zero knowledge inside a VOLE-in-the-Head proof system that
is native to exactly the two ingredients we must prove --- Keccak (a Boolean
circuit) and MAYO verification (a quadratic system over
$\Fsixteen \subseteq \Ftag$). We explain the credential design, the proof
circuit --- including a degree-16 Keccak checkpointing technique that commits
the permutation state only every four rounds --- the quadratic-system
constraint used for MQ verification, the exact 64-bit balance arithmetic, and
the deferred-return settlement extension. We report implementation
performance: on one Apple M4 Pro core-complex a spend proof takes
$\approx$21--23\,ms to produce and $\approx$7.5--9\,ms to verify-and-sign,
with 53\,KiB proofs in the compact profile; tokens themselves are 283 bytes.
\end{abstract}

\section{Introduction}

An anonymous credit token system has one issuer and many clients. The issuer
grants a client a public number of credits; the client can later spend part of
that balance without revealing which issuance --- or which earlier spend ---
produced the token being presented. Each accepted spend consumes a one-time
nullifier (preventing double spending) and returns a fresh token carrying the
remaining balance. The functionality generalizes anonymous tokens in the
Privacy Pass style~\cite{PrivacyPass} from unit tokens to stateful balances,
and specializes keyed-verification anonymous
credentials~\cite{CMZ14} to a single numeric attribute with exact arithmetic.

Classically this is well-trodden ground: blind signatures~\cite{Chaum82},
CL~signatures~\cite{CL01}, algebraic MACs~\cite{CMZ14} and their pairing-based
relatives give compact, efficient constructions. All of them lean on the
algebraic structure of discrete-log or pairing groups, both for the signature
and for the efficient zero-knowledge presentation protocol. Neither survives
a quantum adversary. Worse, the natural repair --- a post-quantum \emph{blind}
signature --- has a discouraging record: lattice-based blind signatures have
been repeatedly broken or weakened, and no candidate has earned the kind of
confidence one wants under real value.

\paragraph{The core idea: sign a digest, hide the preimage.}
VOLE-ACT, following the Katz--Schlesinger design~\cite{KS}, sidesteps
blindness. A credential is a short XOF stream
\[
X \;=\; \shake(\text{domain} \,\|\, \ctx \,\|\, k \,\|\, \enc(b) \,\|\, \rho),
\]
over a client-chosen 256-bit nullifier key $k$, public balance $b$, and
256-bit hiding nonce $\rho$. The first $4m$ bits of $X$ form the
\emph{commitment digest} $\Cred \in \Fsixteen^m$; the next 256 bits form the
\emph{nullifier} $\Null$. At issuance the client \emph{reveals} $\Cred$ and
proves in zero knowledge that it is well-formed for the public balance $b$;
the issuer then computes an ordinary MAYO preimage
$\sigma \gets \Spre(\mathsf{sk}, \Cred)$ and returns it. Nothing is blind:
the issuer sees exactly the digest it signs. Unlinkability comes later --- at
spend time the client proves \emph{in zero knowledge} that it knows
$(\sigma, k, b, \rho)$ with $P(\sigma) = \Cred(k,b,\rho)$, revealing only the
nullifier $\Null$, the public spend amount, and a fresh commitment digest
$\Cred'$ for the successor token. The signature and the signed digest are
never shown again, so there is nothing to link. The unrevealed suffix of $X$
is treated as pseudorandom given its prefix, which is what makes the
nullifier unlinkable to the issuance in which the issuer saw $\Cred$.

This pattern reduces the required primitives to (i) a post-quantum signature
whose \emph{verification} is cheap inside a proof system, and (ii) a
post-quantum zero-knowledge proof system that is efficient on
\emph{hash-heavy Boolean} statements. The pairing of MAYO with
VOLE-in-the-Head is unusually harmonious on both counts:

\begin{itemize}
  \item \textbf{MAYO}~\cite{Beullens21,MAYOspec} verification evaluates a
  public multivariate quadratic map over $\Fsixteen$. Since
  $\Fsixteen$ embeds as a subfield of the proof system's tag field
  $\Ftag$, the whole verification equation becomes a single batched
  \emph{quadratic system} constraint in the proof --- no bit decomposition,
  no modular arithmetic emulation.
  \item \textbf{VOLE-in-the-Head}~\cite{VOLEitH,FAEST} proof systems commit
  witness \emph{bits} and prove Boolean and low-degree relations at
  QuickSilver~\cite{QuickSilver,Wolverine} cost: linear gates are free, and a
  degree-$d$ relation costs one $\Ftag$ element of proof. Keccak --- the only
  hash we use anywhere --- is the friendliest possible circuit for such a
  system: its round function is degree 2.
\end{itemize}

The result is a scheme whose zero-knowledge layer is hash-based (its security
reduces to Keccak's soundness as a correlation-robust hash / random oracle),
whose unforgeability is exactly MAYO's, and whose anonymity survives even if
MAYO falls. There are no lattices, no pairings, and no bespoke assumptions.

\paragraph{Contributions of the implementation described here.}
(1)~A concrete credential and settlement design: one XOF stream yielding both
commitment and nullifier; a typed two-format token state machine; and a
\emph{deferred-return} extension in which the issuer selects a bounded return
$t' \le s$ only \emph{after} verifying the spend proof, wrapped in a nested
hash so the returned amount cannot be reopened or transplanted.
(2)~A proof circuit for SHAKE256 with \emph{degree-16 checkpointing}: the
permutation state is committed only every four rounds, quartering the
dominant witness cost by exploiting the proof system's polynomial constraint
degree.
(3)~A batched quadratic-system constraint that expresses the whipped MAYO
verification equation with shared product terms across all $m$ equations,
folded by verifier randomness.
(4)~An engineering result: the full protocol runs at interactive latencies on
commodity hardware (Table~\ref{tab:perf}) with 283-byte tokens and 53\,KiB
compact spend proofs, using hardware Keccak, word-level bit-matrix
transposition, and parallel GGM tree expansion.

\section{Preliminaries}

\subsection{Notation}
$\lambda = 128$ throughout. $\Ftwo$, $\Fsixteen = \Ftwo[x]/(x^4{+}x{+}1)$ and
$\Ftag = \Ftwo[x]/(x^{128}{+}x^7{+}x^2{+}x{+}1)$ are the Boolean, MAYO, and
tag fields; $\iota : \Fsixteen \hookrightarrow \Ftag$ is the canonical field
embedding (the same reduction polynomial family used by
FAEST~\cite{FAEST}). $\shake$ is used for every PRG, commitment, and
Fiat--Shamir hash, always with prefix-free domain separation and
length-framed inputs~\cite{FIPS202}. $\enc(\cdot)$ is 64-bit little-endian
encoding.

\subsection{VOLE-based zero knowledge and VOLE-in-the-Head}
\label{sec:voleith}

In VOLE-based ZK~\cite{Wolverine,QuickSilver}, the prover holds, for each
witness bit $u_t \in \{0,1\}$, a \emph{tag} $V_t \in \Ftag$, and the verifier
holds a global secret $\Delta \in \Ftag$ and per-coordinate \emph{keys}
\[
K_t \;=\; V_t + u_t \cdot \Delta ,
\]
a vector oblivious linear evaluation (VOLE) correlation. Wires are affine
combinations of committed bits and are linearly homomorphic for free. A
multiplication (degree-2) constraint $a \cdot b = c$ is checked
QuickSilver-style: the verifier forms $B = K_a K_b + \Delta K_c$, which the
prover can match with a degree-1 polynomial in $\Delta$ exactly when the
constraint holds; a random linear combination with challenge weights
$\chi_i$ batches all constraints into one polynomial identity checked at
$\Delta$. The generalization to degree-$d$ polynomial constraints costs
$d$ published coefficients and $d-1$ random VOLE masks in total for the whole
batch.

VOLE-in-the-Head~\cite{VOLEitH} makes this \emph{publicly verifiable and
non-interactive}: the prover generates the VOLE correlation itself from
$\tau$ GGM~\cite{GGM86} seed trees of depth $k$ (with $\tau \cdot k =
\lambda$), commits to all $2^k$ leaves of each tree with an all-but-one
vector commitment, and the Fiat--Shamir challenge $\Delta$ selects, per tree,
one leaf that stays hidden; the remaining leaves are opened and the verifier
recomputes its keys. Soundness comes from the hidden leaf; zero knowledge
from the fact that the opened leaves are exactly what the verifier can
recompute anyway. Our instantiation follows the hardening of
FAEST~\cite{FAEST} and ``One Tree to Rule Them All''~\cite{OneTree}: every
PRG and commitment call binds a per-proof $2\lambda$-bit salt, a per-tree
derived salt, and the exact tree position, eliminating multi-target and
multi-instance batch attacks on the 128-bit seeds; small-field consistency
across the $\tau$ trees uses public corrections re-based onto one shared
$u$-vector, in the style of SoftSpokenOT~\cite{SoftSpoken}.

Three parameter profiles trade proof size against tree-expansion work:
\emph{compact} $(\tau,k) = (16,8)$, \emph{balanced} $(32,4)$, and
\emph{low-latency} $(64,2)$. Each opened proof reveals $\tau$ tree openings
of $k$ sibling seeds, so proofs shrink as $k$ grows while the prover expands
$\tau \cdot 2^k$ leaves.

\paragraph{The consistency check must be a matrix hash.}
Because each tree contributes only $k$ bits of $\Delta$, a cheating prover
who mangles one tree's correction data is caught by a per-tree check only
with probability $1 - 2^{-k}$ --- as low as $3/4$ for $k=2$. The
implementation therefore applies a $128$-row \emph{column-wise linear
universal hash} to the entire tag matrix (two independent polynomial
evaluations combined with random weights, plus a one-time-pad block), so that
any inconsistency survives into the check equation with probability
$1 - 2^{-\lambda}$ regardless of the chunk structure. A scalar hash in place
of the matrix hash is a real soundness bug we caught during development; we
flag it because the failure is silent and parameter-dependent.

\subsection{MAYO}
\label{sec:mayo}

MAYO~\cite{Beullens21,MAYOspec} is an oil-and-vinegar~\cite{UOV} trapdoor
map: a public quadratic map $P^\ast : \Fsixteen^{n} \to \Fsixteen^{m}$
vanishing on a secret $o$-dimensional oil space, with $o < m$ deliberately
too small for direct Kipnis--Shamir-style attacks; the ``whipping''
construction evaluates $k_w$ correlated copies to stretch the oil space, so
the effective public map is
$P : \Fsixteen^{k_w n} \to \Fsixteen^{m}$. The trapdoor samples preimages
$\sigma$ with $P(\sigma) = y$ for arbitrary targets $y$; this
\emph{hash-then-invert} usage (rather than MAYO's message-signing API) is
exactly what the scheme needs. We use the round-2 parameter set
MAYO$_2$: $(n,m,o,k_w) = (81,64,17,4)$, so $\sigma \in \Fsixteen^{324}$
(162 bytes) and targets are $m = 64$ nibbles (32 bytes). We note the risk
posture explicitly: multivariate trapdoors have a difficult history ---
Rainbow fell to a practical attack in 2022~\cite{Beullens22} --- and MAYO's
whipping is comparatively young. The design compartmentalizes this risk:
a MAYO break forges credentials but does \emph{not} deanonymize past
transactions, whose privacy rests only on the hash-based layer.

For the proof circuit, the whipped map is expanded once per public key into
$m$ upper-triangular forms $G_a \in \Fsixteen^{k_w n \times k_w n}$ with
$P(\sigma)_a = \sigma^{\!\top} G_a\, \sigma$; the circuit then proves
$m$ quadratic equations over the \emph{same} $\binom{k_w n + 1}{2}$ product
terms (Section~\ref{sec:quadsys}).

\section{The scheme}
\label{sec:scheme}

\subsection{Credentials and nullifiers from one XOF stream}

Fix a 256-bit context $\ctx$ binding the issuer key (by hash of the expanded
public map), application domain, parameter set, VOLE profile, and protocol
version. The client samples $k, \rho \gets \{0,1\}^{256}$ and, for base
balance $b \in \{0,\dots,2^{64}{-}1\}$, derives the stream
\[
X = \shake(\text{``credential''} \,\|\, \ctx \,\|\, k \,\|\, \enc(b) \,\|\, \rho),
\qquad
\Cred = X[0 : 4m], \quad \Null = X[4m : 4m{+}256].
\]
Both values come from a \emph{single} sponge evaluation --- the message fits
one Keccak rate block --- so proving both costs one in-circuit permutation.
The commitment prefix $\Cred$ is revealed to the issuer exactly once (at
signing time); the nullifier suffix $\Null$ is revealed exactly once (at
spending time). Under the standard modeling of the sponge, the suffix is
pseudorandom given the prefix, which is what divorces the spend from the
issuance.

\subsection{Two token formats and the settlement state machine}

A \emph{direct} token is $(\sigma, k, b, \rho)$ with
$P(\sigma) = \Cred(k,b,\rho)$ and effective balance $b$. A
\emph{deferred-return} token additionally carries an issuer-chosen return
$t$, with
\[
P(\sigma) \;=\; \Def\bigl(\ctx, \Cred(k,b,\rho), t\bigr)
\;=\; \shake\bigl(\text{``deferred''} \,\|\, \ctx \,\|\,
\mathsf{pack16}(\Cred) \,\|\, \enc(t)\bigr)[0:4m],
\]
and effective balance $b + t$ (exact, non-wrapping). The nested hash is
essential: signing any \emph{algebraically adjustable} target (say
$\Cred + \iota(t)$) would let a holder reopen one signature to a different
return. Here changing $t$ requires a fresh preimage or a target collision.

Spends are typed. An ordinary \texttt{spend} consumes either format and
always outputs a direct token; \texttt{spend-with-deferred-return} consumes
either format and outputs a deferred-return token. The client commits to a
fresh $(k', \rho', b')$ \emph{before} the issuer chooses $t' \le s$, which is
precisely the case that cannot be expressed as a smaller net spend --- if the
return were known in advance, the client would simply spend less. All four
transitions carry distinct domain-separation tags end-to-end, so mode
confusion at the byte level cannot create an untyped path.

\subsection{The spend relation}

To spend public amount $s$ from a direct token, the client reveals
$(s, \Null, \Cred')$ and proves knowledge of
$(\sigma, k, b, \rho, k', b', \rho')$ such that
\[
P(\sigma) = \Cred(k,b,\rho), \qquad
\Null = \Null(k,b,\rho), \qquad
b' + s = b, \qquad
\Cred' = \Cred(k',b',\rho').
\]
For a deferred-return input the first equation becomes
$P(\sigma) = \Def(\ctx, \Cred(k,b,\rho), t)$ with hidden $t$, together with
the exact fold $b + t = e$ and $b' + s = e$. The balance equations are
\emph{exact 64-bit unsigned additions with zero final carry}, proven over
committed carry bits --- there is no modular wraparound to exploit, and
overspending is exactly the statement ``the final carry is 1,'' which the
circuit rejects. The issuer verifies, signs the fresh digest
($\sigma' \gets \Spre(\mathsf{sk}, \Cred')$, or the nested target for a
deferred settlement), and atomically stores the request digest and response
under $\Null$; exact retransmissions replay the stored response, so a
network failure cannot burn a balance.

\section{The proof circuit}
\label{sec:circuit}

The relation above is almost entirely SHAKE256: a direct-input spend proves
two hidden sponge evaluations (old stream, fresh stream), a deferred-input
spend proves three (the nested target costs one more). Everything else ---
the MAYO equation and the balance arithmetic --- is quadratic. This section
describes the three techniques that make the circuit small.

\subsection{Degree-16 Keccak checkpointing}
\label{sec:checkpoint}

Keccak-$f[1600]$'s round function is degree 2 (the $\chi$ step is the only
nonlinearity)~\cite{FIPS202}. The naive circuit commits all 1600 state bits
after every round: 24 rounds $\times$ 1600 witness bits, with degree-2
constraints between layers. But witness bits are the dominant cost in a
VOLE-style proof (each consumes a VOLE coordinate, enlarging both the
correction data and every leaf-PRG expansion), while raising the
\emph{degree} of the batched QuickSilver check is cheap: a maximum degree
$d$ costs $d$ published field elements and $d-1$ mask VOLEs \emph{in total}.

We therefore commit the state only every \emph{four} rounds. Between
checkpoints the circuit carries symbolic polynomial expressions: after
rounds $1,2,3,4$ the expressions have degree $2,4,8,16$, and the checkpoint
constraint asserts (per bit) that the degree-16 expression equals the next
committed state --- a degree-16 polynomial constraint, the largest the
implementation accepts. This quarters the per-permutation witness from
$24 \times 1600$ to $6 \times 1600 = 9600$ bits. One permutation dominates
each hidden SHAKE evaluation, so a direct spend's witness is about
$21{,}700$ bits and a deferred spend's about $31{,}500$ --- versus roughly
$79{,}000$ and $118{,}000$ without checkpointing. The prover-side cost is the
schoolbook product of coefficient vectors in the expression layer, which is
pure $\Ftag$ multiplication throughput ($\approx$1.1M multiplications per
permutation) and benefits from carry-less multiply instructions.

\subsection{The quadratic-system constraint for MAYO}
\label{sec:quadsys}

Verifying $P(\sigma) = y$ means checking $m = 64$ quadratic equations
\[
\sum_{r \le c} G_a[r,c]\; \sigma_r \sigma_c \;=\; y_a,
\qquad a \in [m],
\]
over the \emph{same} $\approx \binom{325}{2} \approx 52{,}650$ products
$\sigma_r \sigma_c$. Materializing each product as a committed wire would add
$52{,}650$ witness bits (times four, as these are $\Fsixteen$ values);
instead the backend supports a \emph{deferred quadratic system}: the prover
records the values and tags of every term factor, and at challenge time the
$m$ equations are folded into a single QuickSilver degree-2 check by random
weights $\varphi_a \in \Ftag$ drawn from the transcript, with coefficients
lifted through the embedding $\iota$. An unsatisfied equation survives the
fold except with probability $2^{-\lambda}$. The committed signature nibbles
enter as bit wires ($4 \cdot k_w n = 1296$ bits), and the fold costs the
prover $O(m \cdot |\text{terms}|)$ table-lookup XORs --- data-independent,
since the tables are indexed by public key coefficients. The same mechanism
folds the 64 quadratic carry equations of each balance addition.

\subsection{Balance arithmetic in $\Ftwo$}

Additions are proven bitwise with committed carry vectors:
$a_i + b_i + c_i = s_i$ linearly (over $\Ftwo$, addition is XOR --- an early
design error we caught was writing balance equations additively over the tag
field, where they say nothing about integers), and the carry recurrence
$c_{i+1} = a_i b_i + c_i(a_i + b_i)$ quadratically via the folded system.
Asserting the final carry to be zero makes every balance equation an exact
statement about 64-bit unsigned integers, so conservation is enforced
bit-exactly rather than modulo anything.

\section{Security discussion}
\label{sec:security}

We summarize the intended arguments and their current status; a complete
reduction for the composed system is future work, and the implementation
carries an explicit prototype disclaimer.

\paragraph{Soundness / fiscal safety.}
An accepted spend implies, except with probability negligible in $\lambda$
(via the VOLEitH extraction argument, the matrix consistency check of
Section~\ref{sec:voleith}, and the quadratic-system fold), knowledge of a
MAYO preimage for a well-formed credential digest, correct nullifier
derivation, and exact balance arithmetic. Creating value therefore requires
forging a MAYO preimage or finding structured SHAKE collisions. Double
spending requires two accepted presentations of one credential, which reveal
the same deterministic nullifier and are refused by the nullifier store
(whose atomic retry semantics also make honest retransmission safe).

\paragraph{Anonymity.}
A spend reveals $(s, \Null, \Cred', \text{proof})$. The proof is
zero-knowledge (honest-verifier ZK of the VOLEitH argument, compiled by
Fiat--Shamir with a $2\lambda$-bit per-proof salt); $\Cred'$ is a fresh
commitment whose hiding rests on the 256-bit nonce $\rho'$; and $\Null$ is
the unrevealed suffix of an XOF stream whose prefix the issuer saw at
issuance --- pseudorandom under standard sponge modeling. Consequently
linking a spend to its issuance (or to a prior spend in the same chain)
reduces to breaking the hash layer, \emph{not} the MAYO layer; anonymity
survives a MAYO break. We stress the systemic caveat that spend \emph{amounts}
are public, so deployments have an anonymity set bounded by amount and
context patterns.

\paragraph{Multi-instance hardening.}
Every PRG and commitment call in the tree layer binds salt and position
(after~\cite{OneTree}); a batch adversary attacking $Q$ proofs $\times$ $N$
positions gains no advantage over attacking a single 128-bit seed. The wide
consistency hash removes the $2^{-k}$ per-chunk cheating strategy. Domain
tags separate every hash use; all encodings are canonical, length-framed,
and rejected on any alternate serialization.

\paragraph{Implementation posture.}
All field arithmetic is branch-free with no secret-indexed table lookups
(the $\Fsixteen$ kernels operate on packed lanes with mask-based selection;
$\Ftag$ multiplication uses carry-less multiply instructions); MAYO's
Gaussian elimination runs on a fixed public schedule with masked selection.
Secret-bearing heap state (VOLE tags, expression coefficients, recorded
constraints, tree seeds) is zeroized on drop, including panic paths, and
constraint storage is pre-sized so growth reallocation can never copy
secrets into unwiped freed memory. These are best-effort source-level
properties, documented as such.

\section{Parameters, sizes, and performance}
\label{sec:performance}

All measurements: research prototype, Rust, single Apple M4 Pro (10 cores
used by the parallel tree layer), MAYO$_2$, $\lambda = 128$. Criterion
benchmarks and the measurement harness agree to within a few percent.

\begin{table}[h]
\centering
\caption{Wire sizes (bytes). Tokens are profile-independent; only
proof-carrying requests grow with the low-latency profiles.}
\label{tab:sizes}
\begin{tabular}{lrrr}
\toprule
Object & compact & balanced & low-latency \\
& $(\tau,k){=}(16,8)$ & $(32,4)$ & $(64,2)$ \\
\midrule
Token (either format) & 283 & 283 & 283 \\
Issue request & 29\,750 & 55\,286 & 106\,358 \\
Spend request (direct input) & 52\,990 & 101\,726 & 199\,198 \\
Spend request (deferred input) & 72\,574 & 140\,894 & 277\,534 \\
Issue / spend response & 171 / 179 & 171 / 179 & 171 / 179 \\
Public key (expanded map) & 106\,319 & 106\,319 & 106\,319 \\
\bottomrule
\end{tabular}
\end{table}

\begin{table}[h]
\centering
\caption{Operation latency, median of 15 (ms). ``Verify-and-sign'' is the
complete issuer side including the MAYO preimage computation.}
\label{tab:perf}
\begin{tabular}{lrrr}
\toprule
Operation & compact & balanced & low-latency \\
\midrule
Issue: client prove            & 9.3  & 7.8  & 7.8 \\
Issue: issuer verify-and-sign  & 4.0  & 2.9  & 2.9 \\
Spend (direct): client prove   & 23.0 & 21.5 & 21.6 \\
Spend (direct): issuer verify-and-sign & 9.0 & 7.5 & 7.5 \\
Spend (deferred): client prove & 31.6 & 29.1 & 29.7 \\
Spend (deferred): issuer verify-and-sign & 11.4 & 9.3 & 9.3 \\
Deferred-return: issuer verify-and-sign & 9.0 & 7.5 & 7.5 \\
Issue: end-to-end              & 14.3 & 11.7 & 11.7 \\
Spend (direct): end-to-end     & 33.1 & 30.1 & 30.2 \\
\bottomrule
\end{tabular}
\end{table}

Notably, after optimizing the tree layer (parallel GGM expansion, streaming
pairwise-halving plane accumulation, hardware Keccak via the ARMv8.2-SHA3
extension, and word-level $64{\times}64$ bit-matrix transposition of the
VOLE plane data), the compact profile's latency penalty over low-latency has
collapsed to $\approx$6\% while its proofs are $3.8\times$ smaller --- making
compact the sensible default. An issuer core sustains $>$100
verify-and-sign operations per second; the client prover, running once per
transaction, is comfortably interactive.

\section{Related work}

VOLE-based designated-verifier ZK begins with
Wolverine~\cite{Wolverine} and Mac'n'Cheese~\cite{MacNCheese};
QuickSilver~\cite{QuickSilver} gives the degree-$d$ batched check we
generalize. VOLE-in-the-Head~\cite{VOLEitH} removes the designated verifier
and underlies the FAEST signature candidate~\cite{FAEST}, whose engineering
(GGM trees, salting, consistency checking) we follow, with the hardening of
\cite{OneTree}. MAYO~\cite{Beullens21,MAYOspec} descends from unbalanced
oil-and-vinegar~\cite{UOV}; \cite{Beullens22} is the cautionary tale for the
assumption class. Anonymous tokens and credentials span blind
signatures~\cite{Chaum82}, CL signatures~\cite{CL01}, keyed-verification
credentials from algebraic MACs~\cite{CMZ14}, and deployed unit-token
systems~\cite{PrivacyPass}; VOLE-ACT is a post-quantum stateful-balance
member of this family, instantiating the Katz--Schlesinger design~\cite{KS}.

\section{Open problems}

(1)~A complete security reduction for the composed system (VOLEitH extraction
$+$ quadratic-system fold $+$ the credential/nullifier stream) in the QROM.
(2)~A structured, non-materialized QuickSilver check for the whipped MAYO
equation, removing the $O(m (k_w n)^2)$ expanded-form table and enabling the
larger MAYO parameter sets.
(3)~Hashing the compact public-key encoding (rather than the expanded forms)
into the context --- a transcript-breaking cleanup best done before any
format freeze.
(4)~Aggregating or compressing spend proofs; at 53\,KiB the compact profile
is deployable server-side but remains the dominant bandwidth cost.

\begin{thebibliography}{99}\small

\bibitem{Beullens21}
W.~Beullens.
\emph{MAYO: Practical Post-Quantum Signatures from Oil-and-Vinegar Maps}.
SAC 2021.

\bibitem{Beullens22}
W.~Beullens.
\emph{Breaking Rainbow Takes a Weekend on a Laptop}.
CRYPTO 2022.

\bibitem{MAYOspec}
W.~Beullens, F.~Campos, S.~Celi, B.~Hess, M.~J.~Kannwischer.
\emph{MAYO: Algorithm Specification}.
NIST PQC Additional Digital Signatures, Round 2, 2024.

\bibitem{VOLEitH}
C.~Baum, L.~Braun, C.~Delpech de Saint Guilhem, M.~Kloo{\ss}, E.~Orsini,
L.~Roy, P.~Scholl.
\emph{Publicly Verifiable Zero-Knowledge and Post-Quantum Signatures from
VOLE-in-the-Head}.
CRYPTO 2023.

\bibitem{FAEST}
C.~Baum et al.
\emph{FAEST: Algorithm Specification}.
NIST PQC Additional Digital Signatures, Round 2, 2024.
\url{https://faest.info}.

\bibitem{OneTree}
C.~Baum, W.~Beullens, S.~Mukherjee, E.~Orsini, S.~Ramacher, C.~Rechberger,
L.~Roy, P.~Scholl.
\emph{One Tree to Rule Them All: Optimizing GGM Trees and OWFs for
Post-Quantum Signatures}.
ASIACRYPT 2024.

\bibitem{QuickSilver}
K.~Yang, P.~Sarkar, C.~Weng, X.~Wang.
\emph{QuickSilver: Efficient and Affordable Zero-Knowledge Proofs for
Circuits and Polynomials over Any Field}.
ACM CCS 2021.

\bibitem{Wolverine}
C.~Weng, K.~Yang, J.~Katz, X.~Wang.
\emph{Wolverine: Fast, Scalable, and Communication-Efficient Zero-Knowledge
Proofs for Boolean and Arithmetic Circuits}.
IEEE S\&P 2021.

\bibitem{MacNCheese}
C.~Baum, A.~J.~Malozemoff, M.~B.~Rosen, P.~Scholl.
\emph{Mac'n'Cheese: Zero-Knowledge Proofs for Boolean and Arithmetic
Circuits with Nested Disjunctions}.
CRYPTO 2021.

\bibitem{SoftSpoken}
L.~Roy.
\emph{SoftSpokenOT: Quieter OT Extension from Small-Field Silent VOLE in the
Minicrypt Model}.
CRYPTO 2022.

\bibitem{GGM86}
O.~Goldreich, S.~Goldwasser, S.~Micali.
\emph{How to Construct Random Functions}.
Journal of the ACM 33(4), 1986.

\bibitem{UOV}
A.~Kipnis, J.~Patarin, L.~Goubin.
\emph{Unbalanced Oil and Vinegar Signature Schemes}.
EUROCRYPT 1999.

\bibitem{FIPS202}
NIST.
\emph{SHA-3 Standard: Permutation-Based Hash and Extendable-Output
Functions}.
FIPS PUB 202, 2015.

\bibitem{Chaum82}
D.~Chaum.
\emph{Blind Signatures for Untraceable Payments}.
CRYPTO 1982.

\bibitem{CL01}
J.~Camenisch, A.~Lysyanskaya.
\emph{An Efficient System for Non-transferable Anonymous Credentials with
Optional Anonymity Revocation}.
EUROCRYPT 2001.

\bibitem{CMZ14}
M.~Chase, S.~Meiklejohn, G.~Zaverucha.
\emph{Algebraic MACs and Keyed-Verification Anonymous Credentials}.
ACM CCS 2014.

\bibitem{PrivacyPass}
A.~Davidson, I.~Goldberg, N.~Sullivan, G.~Tankersley, F.~Valsorda.
\emph{Privacy Pass: Bypassing Internet Challenges Anonymously}.
PETS 2018.

\bibitem{KS}
J.~Katz, S.~Schlesinger.
\emph{Anonymous Credit Tokens}.
% NOTE: fill in venue / eprint number for the reference design.

\end{thebibliography}

\end{document}
